\documentclass[11pt]{article}
\usepackage[margin=1in]{geometry}
\usepackage{amsmath}
\usepackage{amssymb}
\usepackage{amsthm}
\usepackage{booktabs}
\usepackage{array}
\usepackage{enumitem}
\usepackage[hidelinks]{hyperref}
\usepackage{xcolor}

\theoremstyle{plain}
\newtheorem{theorem}{Theorem}[section]
\newtheorem{lemma}[theorem]{Lemma}
\newtheorem{corollary}[theorem]{Corollary}
\newtheorem{counterexample}[theorem]{Counterexample}

\newcommand{\lc}{\mathrm{lc}}
\newcommand{\rc}{\mathrm{rc}}
\newcommand{\LCA}{\operatorname{LCA}}
\newcommand{\rev}{\mathrm{rev}}
\newcommand{\R}{\mathbb{R}}

\title{Optimal Leaf Ordering against Trait Distances: a First-Principles Analysis}
\author{}
\date{}

\begin{document}

\maketitle

This document analyses the supplied dynamic program \texttt{olo\_phylo} as a
research candidate. Nothing about the published Bar-Joseph recurrence, the
implementation's state, or the claimed exactness is taken as given. Every claim
is classified as a theorem, a counterexample, or an empirical observation.

\textbf{Intended-use assumptions, stated up front and used only where named.}
The implementation is designed for a \emph{rooted, fully bifurcating} ape
\texttt{phylo} object and a \emph{symmetric} pairwise dissimilarity produced by
\texttt{stats::dist}. Those two hypotheses are \emph{not} part of the
combinatorial problem until they are invoked. They will be shown to be
necessary for exactness of \emph{this} recurrence.

\section{Precise formalization of the optimization problem}

\subsection{Input}

A pair \((T, d)\):

\begin{enumerate}[label=\arabic*.]
\item \textbf{Tree.} A rooted tree \(T = (V, E)\) with leaf set
  \(L(T) = \{1, \dots, n\}\), \(n \ge 2\), in which every internal node has
  \textbf{exactly two children}. Branch lengths are not part of the input.
  Write \(L(v)\) for the leaves in the subtree rooted at \(v\),
  \(\lambda(v) = |L(v)|\), and \(\lc(v),\rc(v)\) for the two children of an
  internal node \(v\). The unique root is \(\rho\). There are exactly \(n-1\)
  internal nodes.
\item \textbf{Dissimilarity.} A real \(n \times n\) matrix \(d\). No metric
  axioms, no non-negativity, and no symmetry are required by the
  \emph{problem}. (The \emph{algorithm} will later impose symmetry.)
\end{enumerate}

The supplied driver builds \(d\) from a trait matrix
\(X \in \R^{n \times p}\) by optional column-standardization and Euclidean
\texttt{dist}. That is one data-generating process for \(d\), not a restriction
on the feasible set.

\subsection{Decision variables}

At each internal node \(v\) a binary \textbf{flip} \(\phi(v) \in \{0,1\}\)
chooses which child is drawn on the left:

\[
\pi_v \;=\;
\begin{cases}
\pi_{\lc(v)}\,\Vert\,\pi_{\rc(v)} & \text{if }\phi(v)=0,\\
\pi_{\rc(v)}\,\Vert\,\pi_{\lc(v)} & \text{if }\phi(v)=1,
\end{cases}
\]

with \(\pi_\ell = (\ell)\) for a leaf \(\ell\). The global layout is
\(\pi = \pi_\rho\). Equivalently, \(\phi\) ranges over \(\{0,1\}^{n-1}\), so
there are \textbf{exactly \(2^{n-1}\)} feasible layouts. They are pairwise
distinct as permutations: flipping \(v\) swaps two nonempty disjoint blocks.

No other moves are permitted. In particular the topology of \(T\) is immutable:
leaves of any subtree must occupy a contiguous block of \(\pi\)
(Lemma~\ref{lem:contig}).

\subsection{Feasibility constraints}

A permutation \(\pi\) of \(L(T)\) is \textbf{\(T\)-consistent} if it arises from
some \(\phi\), iff for every node \(v\) the set \(L(v)\) is a consecutive
interval of \(\pi\). The family \(\{L(v)\}_v\) is laminar, so the
consecutive-block constraints are laminar.

\begin{lemma}[contiguous blocks]\label{lem:contig}
For every \(T\)-consistent \(\pi\) and every node \(v\), the image
\(\pi^{-1}(L(v))\) is an interval of \(\{1,\dots,n\}\).
\end{lemma}

\begin{proof}
Induction on \(\lambda(v)\). A leaf is an interval of length 1. At an internal
node the two child blocks are consecutive by the inductive hypothesis and are
concatenated, hence their union is consecutive.
\end{proof}

\subsection{Objective}

\[
J(\pi)
\;=\;
\sum_{i=1}^{n-1} d\bigl(\pi(i),\,\pi(i+1)\bigr).
\]

There is \textbf{no wrap-around term} \(d(\pi(n),\pi(1))\). The problem is a
\emph{linear} (path) layout, not a circular phylogram layout. The objective
does not depend on branch lengths, on the depth of nodes, or on any
internal-node coordinates.

\subsection{Output, and what ``globally optimal'' means}

A feasible \(\pi^\star\) (equivalently a flip vector \(\phi^\star\)) together
with the scalar \(J(\pi^\star)\), such that

\[
J(\pi^\star)
\;=\;
\min\bigl\{ J(\pi) : \pi\text{ is \(T\)-consistent}\bigr\}.
\]

Ties may be broken arbitrarily. The implementation returns one such order as
tip-index permutation and the scalar cost, then (separately) asks
\texttt{ape::rotateConstr} to realize the order geometrically. The last step is
not part of the combinatorial problem.

\subsection{Structural properties of the problem}

\begin{itemize}
\item \textbf{Finite exact combinatorial optimization} over a set of size
  \(2^{n-1}\).
\item \textbf{Not} unconstrained Hamiltonian path / TSP: the laminar
  consecutive-block constraints are what make a polynomial algorithm possible.
  Unconstrained min-cost Hamiltonian path on \(d\) is NP-hard and has a larger
  feasible set.
\item \textbf{Reversal.} If \(d = d^\top\), then \(J(\rev(\pi)) = J(\pi)\) for
  every \(\pi\), because the undirected adjacent pairs are identical. If \(d\)
  is not symmetric this identity fails, and \(\rev(\pi)\) need not even have a
  simple relation to \(J(\pi)\). Reversal of a \(T\)-consistent order is always
  \(T\)-consistent (flip every internal node).
\item \textbf{No greedy-choice property.} A locally cheapest arrangement of a
  proper subtree need not extend to a global optimum
  (Theorem~\ref{thm:noendpoint}, Counterexample~\ref{ce:5.3}).
\item \textbf{Optimal substructure holds only after the interface is recorded.}
  The interface of a subtree with the rest of the layout is the ordered pair of
  endpoint leaves of its contiguous block (Section~\ref{sec:substructure}).
\item \textbf{Laminar / hierarchical.} The subsets that ever occupy a
  contiguous block are exactly the sets \(L(v)\). This is the difference
  between this DP and Held--Karp TSP (which ranges over all \(2^n\) subsets).
\end{itemize}

\subsection{Implicit assumptions of the implementation (not of the problem)}

The code, as written, additionally assumes:

\begin{center}
\renewcommand{\arraystretch}{1.3}
\begin{tabular}{@{}l p{7.2cm} p{6.2cm}@{}}
\toprule
Assumption & Where it enters & If violated \\
\midrule
\(T\) is rooted & unique node of in-degree 0 & unrooted ape trees are basally
trifurcating and trip the binary check \\
\(T\) is strictly binary & \texttt{length(kids) != 2} is fatal & polytomies
must be resolved \emph{before} the feasible set is defined;
\texttt{multi2di} \emph{changes the problem} \\
\(d\) is symmetric &
\texttt{L\_full[colnames, rownames] <- t(LT)} identifies \(M(u,w)\) with
\(M(w,u)\) & inexact; can report a number that is not the cost of any feasible
order (Counterexample~\ref{ce:3.4}) \\
\(d\) is indexed by tip numbers \texttt{"1"}..\texttt{"n"}, not by
\texttt{tip.label} & \texttt{make\_trait\_dist} relabels after matching &
calling \texttt{olo\_phylo} on a label-indexed matrix silently misaligns \\
finite numeric \(d\) & \texttt{which.min} / \texttt{== min\_cost} &
\texttt{scale()} on a constant trait yields \texttt{NaN} and breaks both the DP
and the argmin \\
\(n \ge 2\) & \texttt{rtree} contract & a one-leaf tree is degenerate \\
\bottomrule
\end{tabular}
\end{center}

\textbf{Most mathematically faithful reading of ambiguities.}

\begin{itemize}
\item The recurrence is the \emph{minimization} dual of Bar-Joseph et al.\
  (2001), who \emph{maximize} a similarity \(S\). With \(d\) in place of \(-S\)
  (or \(c-S\)) the argmin/argmax coincide. The code is a min-cost program, not
  a max-similarity program.
\item ``Left''/``right'' children in ape's edge list are an arbitrary planar
  embedding. The DP searches all embeddings obtainable by child-swaps; it does
  not depend on the input embedding except as a starting point.
\item The root table stores only the orientation ``original-left child on the
  left''. The reverse orientation is implicitly attributed the same cost. This
  is valid whenever the subtree optimum is \textbf{reversal-invariant}.
  Symmetry \(d=d^\top\) is a simple sufficient condition and is the condition
  satisfied by the supplied \texttt{make\_trait\_dist} pipeline. It is not a
  logically necessary condition for every individual instance: some asymmetric
  matrices can happen to induce reversal-invariant optimum tables on a
  particular tree.
\end{itemize}

\section{The existing recurrence, from first principles}
\label{sec:recurrence}

\subsection{States}

For every node \(v\) the implementation stores a rectangular table whose rows
are \(L(\lc(v))\) and whose columns are \(L(\rc(v))\). Write

\[
M_v(u,w)
\qquad
u \in L(\lc(v)),\;
w \in L(\rc(v))
\]

for the stored entry. A dense ``full'' matrix \(F_v\) over \(L(v) \times L(v)\)
is then materialized by

\[
F_v(u,w)
\;=\;
\begin{cases}
0 & \text{if }\lambda(v)=1,\\
M_v(u,w) & \text{if }(u,w)\text{ is a stored pair},\\
M_v(w,u) & \text{if }(w,u)\text{ is a stored pair (the transpose fill)},\\
+\infty & \text{otherwise}.
\end{cases}
\]

\textbf{Intended meaning of \(M_v(u,w)\).} The minimum of \(J\) restricted to
the contiguous block of \(L(v)\), over \(T\)-consistent arrangements of that
block that start at \(u\) and end at \(w\), \emph{with \(\lc(v)\) placed to the
left of \(\rc(v)\)}.

\textbf{Intended meaning of \(F_v(u,w)\).} The same quantity with \emph{either}
orientation of \(v\)'s two children allowed. The identification
\(F_v(w,u) = F_v(u,w)\) is the reversal shortcut.

\subsection{State variables, listed}

\begin{center}
\renewcommand{\arraystretch}{1.3}
\begin{tabular}{@{}l l p{9.5cm}@{}}
\toprule
Symbol & Domain & Role \\
\midrule
\(v\) & nodes of \(T\) & which laminar block is being arranged \\
\(u\) & a leaf of the left child & left endpoint of the block \\
\(w\) & a leaf of the right child & right endpoint of the block \\
\(M_v(u,w) \in \R\cup\{+\infty\}\) & scalar & best internal path cost for that
interface \\
\(B_v(u,w),\,C_v(u,w)\) & leaves & argmin seam witnesses, for traceback \\
\bottomrule
\end{tabular}
\end{center}

No other history (flip vector, internal order, ancestor context) is stored.

\subsection{Transition}

Let \(L = \lc(v)\), \(R = \rc(v)\). The code computes, for every
\(u \in L(L)\) and \(w \in L(R)\),

\begin{equation}
\label{eq:R}
M_v(u,w)
\;=\;
\min_{\substack{b \in L(L)\\ c \in L(R)}}
\Bigl(
F_L(u,b) \;+\; F_R(c,w) \;+\; d(b,c)
\Bigr),
\end{equation}

skipping pairs \((u,b)\) for which \(F_L(u,b) = +\infty\) (equivalently:
skipping pairs that are not a valid endpoint pair of \(L\)). The inner minimum
over \(c\) is vectorized; the outer minimum over \(b\) is an R loop. Witnesses
\((b^\star, c^\star)\) are stored as \texttt{ptrB}, \texttt{ptrC}.

\subsection{Base case}

A leaf \(\ell\): \(L(\ell) = \{\ell\}\), \(F_\ell(\ell,\ell) = 0\), and there is
no table \(M\). The empty path has cost 0.

\subsection{Reconstruction}

At the root, \(\mathrm{cost} = \min M_\rho\), and a minimizing stored pair
\((u^\star, w^\star)\) is selected (the first, in the event of ties).
Traceback:

\begin{itemize}
\item if the requested \((s,t)\) is a stored (left-start, right-end) pair,
  concatenate \texttt{get\_order(L, s, B(s,t))} with
  \texttt{get\_order(R, C(s,t), t)};
\item if \((s,t)\) is the reverse of a stored pair, return the reverse of
  \texttt{get\_order(v, t, s)};
\item a leaf returns itself.
\end{itemize}

The reconstructed sequence is claimed to be a \(T\)-consistent order of cost
equal to the reported minimum.

\subsection{Does (R) consider every feasible solution?}

\textbf{Every \(T\)-consistent order of \(L(v)\) with left child on the left
and endpoints \((u,w)\)} has a unique seam \((b,c)\): \(b\) is the last leaf of
the left block, \(c\) the first leaf of the right block. The min in (R)
therefore ranges over a set of pairs that \emph{covers} every such order. It
does \textbf{not} enumerate every internal flip configuration: for each
candidate seam it substitutes the already-minimized \(F_L(u,b)\) and
\(F_R(c,w)\). Whether that substitution preserves the global min is exactly the
optimal-substructure question (Section~\ref{sec:substructure}). It is
\emph{not} granted by the existence of a nested min.

Orders with the right child on the left are \textbf{not} stored in \(M_v\).
They are represented only by the transpose fill of \(F_v\). That representation
is correct exactly when the relevant subtree optimum table is
reversal-invariant, that is, when the true endpoint-conditioned optimum
satisfies \(C_v(x,y)=C_v(y,x)\) for every feasible endpoint pair. Symmetry
\(d=d^\top\) implies this property for every subtree, but it is stronger than
necessary on a fixed instance.

\subsection{Are all transitions valid?}

A finite sum \(F_L(u,b)+F_R(c,w)+d(b,c)\) corresponds to concatenating a
\(T\)-consistent order of \(L\) with endpoints \((u,b)\) and a \(T\)-consistent
order of \(R\) with endpoints \((c,w)\). The concatenation is \(T\)-consistent
for \(v\) with \(\phi(v)=0\). No infeasible layout is generated by a
\emph{finite} transition.

The \(+\infty\) entries of \(F\) are exactly the pairs of leaves that lie in
the \emph{same} child of the node (or the diagonal of a non-leaf). Those pairs
cannot be the two endpoints of a binary node (Lemma~\ref{lem:crosschild}), so
excluding them does not drop a feasible layout.

\begin{lemma}[cross-child endpoints]\label{lem:crosschild}
If \(\lambda(v) \ge 2\), then in every \(T\)-consistent order of \(L(v)\) the
two endpoints lie in different children of \(v\). Every pair
\((u,w) \in L(\lc(v)) \times L(\rc(v))\) \emph{is} realized by some
\(T\)-consistent order of \(L(v)\).
\end{lemma}

\begin{proof}
The two children occupy the two sides of the concatenation, so the left
endpoint is in the left-placed child and the right endpoint in the right-placed
child. Conversely, any leaf of a subtree can be made its left endpoint (place
the unique child containing it on the left, recurse) and, independently, any
leaf of the other child can be made the right endpoint. Independence holds
because the two children's internal flips do not interact.
\end{proof}

\subsection{Multiple representations of the same feasible solution}

A single global order \(\pi\) induces, at each \(v\), a unique pair of
endpoints of the block \(L(v)\) and a unique orientation of \(v\)'s children.
So a given \(\pi\) occupies \textbf{exactly one} stored cell of \(M_v\)
\emph{or} its transpose image in \(F_v\), not both as independently optimized
objects. Under \(d = d^\top\) those two cells are forced equal, so
double-counting is only a symmetry, not an alternate cost. Under asymmetric
\(d\) the forced equality is a \emph{false identification of two different
layouts}.

This does not by itself cause overcounting of the reported min (the root min is
taken over stored cells only), but it \emph{does} cause the stored number to be
the cost of a different orientation than the one reconstructed, whenever the
transpose fill was used as an input to (R). That is an invalid transition under
asymmetric \(d\).

\subsection{Can (R) exclude a globally optimal solution?}

Two exclusion mechanisms exist.

\begin{enumerate}[label=\arabic*.]
\item \textbf{Missing orientation at a node, compensated by transpose.}
  Harmless iff \(d = d^\top\); fatal otherwise (Section~\ref{sec:exact}).
\item \textbf{Replacing a sub-arrangement by a cheaper one with the same
  endpoints.} Harmless \emph{if} the state \((v,u,w)\) is sufficient
  (Section~\ref{sec:state}). Harmful if the state forgets an endpoint
  (Section~\ref{sec:state}, \ref{sec:substructure}).
\end{enumerate}

No other pruning is performed. In particular the code does not implement the
Bar-Joseph similarity-bound early exit.

\subsection{Can (R) generate an infeasible solution?}

Traceback concatenates child orders along witnesses of (R). By
Section~\ref{sec:recurrence} every finite witness is a feasible concatenation.
The reconstructed permutation is therefore \(T\)-consistent,
\textbf{provided} the witnesses were computed from correct child tables. Under
asymmetric \(d\) the child tables themselves contain transpose-fill values that
do not equal any feasible child-path cost; traceback still emits a feasible
permutation, but the \emph{number stored in the table need not equal that
permutation's true \(J\)} (Counterexample~\ref{ce:3.4}).

\subsection{Can the reported value beat the true optimum?}

Yes, but only through the invalid transpose identification. Under symmetric
\(d\), every stored number is a feasible-path cost
(Theorem~\ref{thm:exactness}), hence cannot undershoot. Under asymmetric \(d\),
the stored min can be strictly smaller than \(\min\{J(\pi):\pi\text{
\(T\)-consistent}\}\) and need not equal \(J\) of the reconstructed order
(Counterexample~\ref{ce:3.4}, the \(n=4\) directed instance).

\section{Global exactness}
\label{sec:exact}

\subsection{Statement}

\begin{theorem}[exactness under symmetry]\label{thm:exactness}
Assume \(T\) is rooted and strictly binary, and \(d = d^\top\) with entries in
\(\R\) (no \(\mathrm{NaN}\)). Then \texttt{olo\_phylo} returns a
\(T\)-consistent order \(\pi\) with \(J(\pi)\) equal to the global minimum of
\(J\) over all \(T\)-consistent orders.
\end{theorem}

\begin{theorem}[no unconditional guarantee for asymmetric input]\label{thm:asym}
If arbitrary asymmetric \(d\) is allowed, the algorithm is not globally exact
as an algorithmic guarantee. It can (i) return a suboptimal feasible order and,
on larger instances, (ii) report a scalar that is not the \(J\)-value of the
reconstructed order because transpose-filled child costs can be fictitious.
The smallest counterexample to global optimality already has \(n=2\).
\end{theorem}

More precisely, symmetry of \(d\) is \textbf{sufficient but not instance-wise
necessary}. The weakest clean state-level condition is reversal invariance of
every true subtree endpoint table \(C_v\):
\[
C_v(x,y)=C_v(y,x)
\]
for every internal \(v\) and every feasible endpoint pair. Under this condition
the transpose fill is valid and the proof of
Theorem~\ref{thm:exactness} goes through verbatim. Outside this condition the
code may still be accidentally correct on a particular instance, but no general
guarantee remains.

\begin{theorem}[binary-domain restriction]\label{thm:binary}
The implementation requires every internal node to have \textbf{outdegree
exactly 2} in the rooted \texttt{phylo} representation. If a node has any other
number of children, it stops. This is a domain restriction of recurrence (R):
the transition is written for the concatenation of exactly two child blocks.
A multifurcating node with \(k\) children introduces \(k!\) possible child
orders (before accounting for descendants), so the feasible family and the
transition must be generalized rather than silently resolved. Calling
\texttt{multi2di()} first changes that feasible family by choosing a binary
refinement.
\end{theorem}

The rest of this section proves~\ref{thm:exactness} and~\ref{thm:asym}.
Non-negativity of \(d\) is \textbf{not} used.

\subsection{Invariant}

\textbf{Invariant \(I(v)\).} After \texttt{dp(v)} returns,

\begin{enumerate}[label=\arabic*.]
\item For every \(u \in L(\lc(v))\), \(w \in L(\rc(v))\), \(M_v(u,w)\) equals
  the minimum internal path cost of a \(T\)-consistent arrangement of \(L(v)\)
  that starts at \(u\), ends at \(w\), and places \(\lc(v)\) on the left.
\item \(F_v(x,y)\) equals the minimum internal path cost of a \(T\)-consistent
  arrangement of \(L(v)\) that starts at \(x\) and ends at \(y\), over
  \emph{both} child orientations --- \textbf{provided} \(d = d^\top\). If
  \(x,y\) is not a cross-child pair, \(F_v(x,y) = +\infty\).
\item The witnesses \(B_v, C_v\) realize the min in (R).
\end{enumerate}

\subsection{Proof of Theorem~\ref{thm:exactness}}

The argument is induction on \(\lambda(v)\), then a root-finalization step.
Throughout, \(d = d^\top\).

\textbf{Base.} \(\lambda(v)=1\): \(I(v)\) holds (\(F=0\) on the unique pair).

\textbf{Inductive step.} Let \(L,R\) be the children of \(v\), and assume
\(I(L)\) and \(I(R)\). Fix \(u \in L(L)\), \(w \in L(R)\).

\emph{Soundness (\(\ge\)).} Let \(\sigma\) be any \(T\)-consistent arrangement
of \(L(v)\) with \(\lc(v)\) on the left, start \(u\), end \(w\). Let \((b,c)\)
be its seam. The left block is a \(T\)-consistent arrangement of \(L(L)\) with
endpoints \((u,b)\); by \(I(L)\) its internal cost is at least \(F_L(u,b)\).
Likewise the right block costs at least \(F_R(c,w)\). The unique cross edge
contributes \(d(b,c)\). Hence \(J(\sigma) \ge\) the right-hand side of (R), so
\(J(\sigma) \ge M_v(u,w)\).

\emph{Achievability (\(\le\)).} Let \((b^\star,c^\star)\) be a minimizer of (R).
By \(I(L)\) there is a \(T\)-consistent arrangement of \(L(L)\) with endpoints
\((u,b^\star)\) of cost \(F_L(u,b^\star)\); similarly for \(R\). Concatenating
them is \(T\)-consistent for \(v\), starts at \(u\), ends at \(w\), and has cost
exactly the right-hand side of (R). Lemma~\ref{lem:crosschild} guarantees that
every stored pair \((u,w)\) has at least one finite seam (e.g.\ any valid
endpoints of the two children), so \(M_v(u,w) < +\infty\).

Thus (1) holds at \(v\).

\emph{Transpose fill, using symmetry.} Take a \(T\)-consistent arrangement
\(\sigma\) of \(L(v)\) with start \(x\) and end \(y\). Its reverse
\(\rev(\sigma)\) is \(T\)-consistent, starts at \(y\), ends at \(x\), and has
\(J(\rev(\sigma)) = J(\sigma)\) because \(d = d^\top\). Therefore the min cost
with endpoints \((x,y)\) equals the min cost with endpoints \((y,x)\).
Combined with (1) at \(v\) and at the two possible orientations, this is (2).

Witnesses (3) are stored by construction.

\textbf{Finalization at the root.} The stored cells of \(M_\rho\) are exactly
the layouts with \(\lc(\rho)\) on the left. Every \(T\)-consistent global order
either has that orientation or is the reverse of one that does. Under
\(d = d^\top\) the two classes have identical cost sets, so
\(\min M_\rho = \min J\). Traceback, by (3) and the inductive construction,
emits an order realizing that cell.

\textbf{Completeness, soundness, optimality, termination --- mapped to the
requested skeleton.}

\begin{enumerate}[label=\arabic*.]
\item \textbf{Completeness.} Every feasible \(\pi\) induces, at each \(v\),
  some endpoint pair and orientation. That pair is a stored cell or its
  reverse. The min in (R) ranges over all seams of that cell. Induction
  therefore represents every feasible \(\pi\) by a path of states.
\item \textbf{Soundness.} Every finite state/transition is a concatenation of
  \(T\)-consistent child orders (Section~\ref{sec:recurrence}).
\item \textbf{Optimality of retention.} For a fixed \((v,u,w)\), only the min
  cost is kept. The exchange argument of Section~\ref{sec:state} shows that
  discarding worse same-endpoint arrangements cannot drop a global optimum.
\item \textbf{Termination.} The recursion follows the tree, which is finite
  and acyclic. The reported root cell is the global min by the finalization
  paragraph.
\end{enumerate}

\subsection{Proof of Theorem~\ref{thm:asym}, with the smallest counterexample}

\begin{counterexample}[minimal, \(n=2\), asymmetric \(d\)]\label{ce:3.4}
Let the tree be the single cherry \(T=(1,2)\), and let
\[
d(1,2)=5, \qquad d(2,1)=1.
\]
There are exactly two feasible orders:

\begin{center}
\begin{tabular}{lc}
\toprule
order & true cost \\
\midrule
\(1,2\) & \(5\) \\
\(2,1\) & \(1\) \\
\bottomrule
\end{tabular}
\end{center}

The implementation stores only the root cell \(M_\rho(1,2)=5\), takes its
minimum, and returns order \(1,2\) with cost \(5\). The global optimum is the
reverse order \(2,1\) with cost \(1\). Thus \(n=2\) is already sufficient to
refute exactness for arbitrary asymmetric matrices.

This example isolates a \textbf{root-finalization} failure: the code searches
one root orientation and relies on reversal to cover the other. No child
transpose is even needed.
\end{counterexample}

\textbf{Ghost-cost failure requires a larger tree.} Once an internal child
table is used through the transpose fill, the stored scalar can cease to be the
cost of the path later reconstructed. For example, on \(T=((1,2),(3,4))\),
choose asymmetric entries so that \(d(2,1)\neq d(1,2)\) and/or
\(d(4,3)\neq d(3,4)\). The child table stores the forward value but the parent
may consume it as though it were the reverse value. A parent transition can
therefore add costs belonging to mutually inconsistent orientations. The
traceback still emits a topology-feasible order, but the reported scalar need
not equal that order's true directed path cost.

This distinction matters:

\begin{enumerate}[label=\arabic*.]
\item \(n=2\) proves \textbf{suboptimality} under asymmetric \(d\).
\item \(n\ge 4\) can exhibit the stronger \textbf{cost/order inconsistency}
  caused by a transpose-filled child state.
\end{enumerate}

The defect is entirely the unproved reversal identification; storing both
ordered endpoint orientations removes it while preserving the same DP logic.

\subsection{What the intended pipeline actually satisfies}

\texttt{make\_trait\_dist} calls \texttt{dist()}, which returns a symmetric
matrix. For every input that follows the documented driver (\texttt{scale} +
Euclidean / Minkowski \texttt{dist} on real traits without \texttt{NaN}),
Theorem~\ref{thm:exactness} applies and the algorithm is globally exact.

Exactness is \textbf{conditional}, not universal. Symmetry \(d=d^\top\) is the
simple, intended sufficient condition, not ``\(d\) comes from traits''. A
caller who passes an asymmetric Kullback--Leibler matrix, a directed hit-and-miss
score, or a raw non-symmetric matrix leaves the guaranteed-correct domain. Such
an instance may happen to be solved correctly, but the algorithm no longer has
a correctness guarantee (Counterexample~\ref{ce:3.4}).

\section{Minimal sufficient state}
\label{sec:state}

\subsection{The candidate and the true state index}

After processing a subtree block, every objective term outside that block can
interact with its internal arrangement only through the two leaves exposed at
its ends. Thus the natural interface is the endpoint pair.

For a fixed tree, the node identifier is actually \textbf{derivable} from the
two endpoints. In every feasible non-leaf state the endpoints lie in different
children of their subtree root, hence
\[
v=\LCA_T(u,w).
\]
Therefore the mathematically minimal generic state \emph{index} for the
standard subtree-composition DP is:

\begin{itemize}
\item under symmetric \(d\): the unordered endpoint pair \(\{u,w\}\);
\item under directed/asymmetric \(d\): the ordered endpoint pair \((u,w)\).
\end{itemize}

The DP value \(F_v(u,w)\) is the \textbf{value stored at the state}, not an
additional decision-state coordinate. The implementation organizes these pair
states by their derived node \(v\), which is convenient but logically
redundant.

Equivalently, one may keep the notation \((v,u,w)\) for readability, with the
invariant \(v=\LCA(u,w)\). The proof below is about sufficiency of the endpoint
interface, not about a particular hash key or array layout.

\subsection{Classification of each component}

\begin{center}
\renewcommand{\arraystretch}{1.35}
\begin{tabular}{@{}l l p{7.6cm}@{}}
\toprule
Component & Classification & Reason \\
\midrule
Node \(v\) & \textbf{Derivable} from \(u,w\) & any feasible endpoint pair for
a non-leaf state satisfies \(v=\LCA(u,w)\) \\
Endpoint identity \(u\) & \textbf{Necessary} in a generic compositional
summary & an external adjacency can contribute \(d(x,u)\) \\
Endpoint identity \(w\) & \textbf{Necessary} & symmetrically, and it may form
the sibling seam \\
Order of the pair & \textbf{Redundant under symmetric \(d\)}; \textbf{necessary
in general} & reversal preserves the value only in the reversal-invariant case
\\
DP value \(F_v(u,w)\) & \textbf{Payload, not a state coordinate} & it is the
optimal value associated with the endpoint state \\
Full internal order & \textbf{Redundant for value optimization} once endpoints
and optimal value are known & Theorem~\ref{thm:suff} \\
Flip vector inside \(v\) & \textbf{Redundant} for future costs & same exchange
argument \\
Ancestor context & \textbf{Redundant} for this subtree summary & the outside
sees only the two boundary leaves \\
Same-child endpoint pairs & \textbf{Infeasible} and removable &
Lemma~\ref{lem:crosschild} \\
\bottomrule
\end{tabular}
\end{center}

This establishes minimality \textbf{within the natural compositional model
that summarizes an already-solved subtree for arbitrary future parent
contexts}. It does not imply that every exact algorithm must materialize all
endpoint states simultaneously. Section~\ref{sec:algo} discusses an
asymptotically optimal \(O(n^2)\) algorithm using a different sweeping
decomposition and only linear additional working memory.

\subsection{Proof that \((v,u,w)\) is sufficient}

\begin{theorem}\label{thm:suff}
Let \(\sigma, \sigma'\) be two \(T\)-consistent arrangements of \(L(v)\) with
the same endpoints \((u,w)\). For every \(T\)-consistent global order \(\pi\)
whose restriction to \(L(v)\) is \(\sigma\), the order \(\pi'\) obtained by
substituting \(\sigma'\) for that interval is \(T\)-consistent, and
\[
J(\pi') - J(\pi)
\;=\;
J(\sigma') - J(\sigma).
\]
In particular, if \(J(\sigma') \le J(\sigma)\), then \(\pi'\) is at least as
good as \(\pi\). Therefore any two partial solutions with the same
\((v,u,w)\) have \textbf{identical optimal continuation behaviour}: the
continuation sees only \(u\), \(w\), and the already-paid internal cost.
\end{theorem}

\begin{proof}
Substituting \(\sigma'\) does not change which set occupies which interval
(still \(L(v)\)), and \(\sigma'\) is \(T\)-consistent for \(v\), so \(\pi'\)
is \(T\)-consistent. The adjacent pairs of \(\pi\) that are not internal to the
block are: the pair immediately left of \(u\) (if any) and the pair immediately
right of \(w\) (if any). Both involve \(u\) or \(w\) and a leaf outside
\(L(v)\), hence are identical for \(\pi\) and \(\pi'\). All other external pairs
are unchanged. The difference in \(J\) is therefore exactly the difference of
internal path costs.
\end{proof}

\subsection{Proof that neither endpoint can be dropped}

\begin{theorem}[both endpoint identities are necessary in the generic exact
state]\label{thm:twoend}
Even for symmetric \(d\), a state that retains only one endpoint of a subtree
is not sufficient for exact continuation. Under asymmetric \(d\), even an
unordered pair \(\{u,w\}\) is insufficient because orientation matters.
\end{theorem}

\begin{proof}[Proof, dropping the right endpoint while retaining the left]
Consider the subtree
\[
S=((1,2),3)
\]
inside the larger tree \(T=(S,4)\). Use symmetric distances with
\[
d(3,1)=0,\quad d(1,2)=0,\quad d(3,2)=1,
\]
\[
d(2,4)=100,\quad d(1,4)=0,
\]
and choose all other unspecified distances nonnegative and large enough that
they do not create a cheaper competing case for the particular parent state
considered below.

Two feasible arrangements of \(S\) have the \textbf{same left endpoint} \(3\):
\[
\sigma_A=(3,1,2), \qquad J(\sigma_A)=0,
\]
\[
\sigma_B=(3,2,1), \qquad J(\sigma_B)=1.
\]
Their right endpoints are respectively \(2\) and \(1\). If a compressed state
indexed only by \((S,\text{left}=3)\) keeps the cheaper internal layout, it
keeps \(\sigma_A\) and discards \(\sigma_B\). But at the parent, with leaf
\(4\) placed to the right,
\[
J(\sigma_A\Vert 4)=0+d(2,4)=100,
\]
whereas
\[
J(\sigma_B\Vert 4)=1+d(1,4)=1.
\]
Thus the internally worse partial solution becomes strictly better after
continuation. The forgotten right endpoint changes the seam cost and is
therefore necessary.

Dropping the left endpoint is symmetric: reverse the construction (or attach
the external leaf on the left). Hence neither endpoint identity can be removed
from a generic exact state.
\end{proof}

\textbf{Unordered pair \(\{u,w\}\).} Under symmetric \(d\),
\(F_v(u,w)=F_v(w,u)\), so an unordered pair plus a single cost is sufficient if
the combination logic is allowed to choose which endpoint is exposed on each
side. The implementation realizes the same reduction more concretely by storing
one cross-child orientation and obtaining the reverse by \texttt{rev()}. Under
asymmetric \(d\), the ordered states \((u,w)\) and \((w,u)\) are not equivalent
in general.

\subsection{Are all cross-child pairs necessary?}

\begin{theorem}[no generic pair can be dropped]\label{thm:nodrop}
For a node \(v\) with children \(L,R\), and for generic \(d\), every pair
\((u,w) \in L(L)\times L(R)\) is the unique minimizing interface of \(L(v)\)
for some parent context. Hence the \(\lambda(L)\cdot\lambda(R)\) stored cells
are necessary in the worst case over \(d\).
\end{theorem}

\begin{proof}[Proof sketch]
Give \(u\) a unique external left neighbour \(x\) with \(d(x,u) = 0\) and
\(d(x,\cdot) = M\) huge otherwise; give \(w\) a unique external right neighbour
\(y\) similarly; set internal distances so that every arrangement of \(L(v)\)
has comparable internal cost. The unique optimum of the three-block instance
\(x \vert L(v) \vert y\) uses endpoints \((u,w)\). A DP that had dropped that
cell would be forced to use a costly external adjacency.
\end{proof}

This rules out ``keep only \(O(\lambda(v))\) endpoints per node'' as a
\emph{uniform} exact reduction. Special trees (combs: Section~\ref{sec:comb})
\emph{do} have fewer feasible pairs, and a DP that iterates only finite
\(F\)-entries automatically exploits that; the current nested loops do not.


\section{Optimal substructure}
\label{sec:substructure}

Analyse the \emph{problem}, not the code.

\subsection{What holds}

\begin{theorem}[optimal substructure for the endpoint state]\label{thm:optend}
Let \(\pi^\star\) be a globally optimal \(T\)-consistent order, \(v\) any node,
and \((u,w)\) the endpoints of the interval \(\pi^\star\) occupies on \(L(v)\).
Then the restriction \(\pi^\star|_{L(v)}\) is an optimal arrangement of \(L(v)\)
among those with endpoints \((u,w)\).
\end{theorem}

\begin{proof}
Immediate from Theorem~\ref{thm:suff}: a cheaper same-endpoint restriction
could be substituted, preserving feasibility and strictly decreasing \(J\),
contradicting global optimality of \(\pi^\star\).
\end{proof}

This is the Bellman property that (R) relies on. It uses:

\begin{itemize}
\item additivity of \(J\) along the path,
\item Lemma~\ref{lem:contig} (no edges skip over the block),
\item the state containing \emph{both} endpoints.
\end{itemize}

It does \textbf{not} use non-negativity, the metric axioms, or symmetry.
Symmetry is used only to share \(F(u,w)\) with \(F(w,u)\).

\textbf{Overlapping subproblems.} At the endpoint-state level, yes: a child
value \(F_v(u,w)\) is reused by many parent seam candidates. At the
\emph{node-recursion} level, however, the rooted tree has unique parents, so
\texttt{memo[[node]]} is not needed to prevent repeated recursive visits when
starting once at the root; it merely caches the already-built table. The true
state-dependency graph has vertices equal to endpoint-conditioned states and
hyperarcs corresponding to seam combinations. It is acyclic because every
dependency goes from a subtree to a strict descendant subtree.

\textbf{Markov / sufficient-state structure.} Theorem~\ref{thm:suff} is exactly
that: the future is conditionally independent of the internal arrangement given
\((v,u,w)\) and the paid cost.

\textbf{Greedy-choice property.} \textbf{False.} See
Sections~\ref{sec:fails}--\ref{sec:minimalce}. Consequently a ``flip each node
to minimize the present seam, given already-chosen child orders'' algorithm is
not exact, even though it inspects a DP-like tree.

\subsection{What fails: endpoint-free ``optimal subtrees''}
\label{sec:fails}

The statement

\begin{quote}
every optimal solution restricts to an \emph{unconstrained} optimal solution
of each subtree
\end{quote}

is \textbf{false}. The restriction is optimal only \emph{in the slice of fixed
endpoints}.

\begin{theorem}\label{thm:noendpoint}
The problem does \textbf{not} have optimal substructure with respect to the
state ``subtree \(v\)'' (a single best order per node).
\end{theorem}

\subsection{Minimal counterexample (also the necessity witness for
endpoints)}
\label{sec:minimalce}

\begin{counterexample}\label{ce:5.3}
Tree \(T = \bigl(1,\,((2,3),4)\bigr)\). Symmetric distances:
\[
\begin{aligned}
d(2,3) &= 1, &
d(3,4) &= 1, &
d(2,4) &= 0.5,\\
d(1,2) &= 0, &
d(1,3) &= d(1,4) = 100.
\end{aligned}
\]
(All unspecified pairs are irrelevant because they never appear as
\(T\)-adjacencies.)

Let \(T_0 = ((2,3),4)\). The \(T_0\)-consistent orders and internal costs:

\begin{center}
\begin{tabular}{lcc}
\toprule
order of \(T_0\) & internal \(J\) & endpoints \\
\midrule
\(3,2,4\) & \(1 + 0.5 = 1.5\) & \(3,4\) \\
\(4,2,3\) & \(0.5 + 1 = 1.5\) & \(4,3\) \\
\(2,3,4\) & \(1 + 1 = 2\) & \(2,4\) \\
\(4,3,2\) & \(1 + 1 = 2\) & \(4,2\) \\
\bottomrule
\end{tabular}
\end{center}

An algorithm that keeps one best order of \(T_0\) keeps a cost-\(1.5\) layout,
with endpoints in \(\{3,4\}\). Completing with leaf \(1\):

\begin{itemize}
\item \(1 \vert 3,2,4\) costs \(100 + 1.5 = 101.5\),
\item \(1 \vert 4,2,3\) costs \(100 + 1.5 = 101.5\).
\end{itemize}

The globally optimal order is \(1,2,3,4\) with \(J = 0+1+1 = 2\), whose
restriction to \(T_0\) is the \emph{locally suboptimal} layout \(2,3,4\).

The endpoint-aware DP computes \(F_{T_0}(2,4) = 2\) and \(F_{T_0}(3,4) = 1.5\),
then at the root
\[
M_\rho(1,4)
\;=\;
\min_c \bigl(F_{T_0}(c,4) + d(1,c)\bigr)
\;=\;
\min\bigl(2+0,\; 1.5+100\bigr)
\;=\;
2,
\]
and recovers \(1,2,3,4\). Independently, a ``single best order per subtree''
implementation on this instance returns cost \(101.5\).

This counterexample is minimal: \(n=4\), and on \(n=3\) every subtree is a
cherry (internal cost has no alternative), so endpoint-free search happens to
be exact for \(n \le 3\).
\end{counterexample}

\textbf{What additional state restores the property.} Exactly the pair of
endpoints, i.e.\ the state already used by (R) under
Theorem~\ref{thm:optend}.

\section{Dominance relations}
\label{sec:dom}

A dominance rule is safe only if every global optimum that used the discarded
state has a matching (or better) global optimum that uses the retained state.

\subsection{Rules that \emph{are} safe}

\textbf{D1. Same-interface min (already in the DP).} For fixed \((v,u,w)\),
only the smallest internal cost is kept. Safe by Theorem~\ref{thm:suff}. This
is scalar dominance on a singleton interface.

\textbf{D2. Reversal identification, \(d = d^\top\).} The state \((v,u,w)\) and
\((v,w,u)\) are equivalent. Safe by reversal invariance. This is a symmetry
reduction of factor exactly \(2\) on non-leaf tables, and is already implemented.

\textbf{D3. Infeasible same-child pairs.} \(F_v(x,y) = +\infty\) when \(x,y\)
lie in the same child. Safe by Lemma~\ref{lem:crosschild}. Already implemented
(those cells are never written, and are skipped by the \texttt{is.finite}
guard).

\textbf{D4. Pointwise seam dominance (combination-time).} Fix a node \(v\)
with children \(L,R\), and fix a left endpoint \(u \in L(L)\). If
\(b, b' \in L(L)\) satisfy
\[
F_L(u,b) \le F_L(u,b')
\qquad\text{and}\qquad
d(b,c) \le d(b',c)
\quad\text{for all }c \in L(R),
\]
then \(b\) dominates \(b'\) as a seam candidate for this \(u\) against this
sibling \(R\): every continuation that would have used \(b'\) can use \(b\)
instead, and the objective does not increase. \emph{Proof.} The combination
formula (R) adds \(F_L(u,\cdot) + d(\cdot,c) + F_R(c,w)\); the first two terms
are never worse for \(b\) than for \(b'\), for every \(c\), and \(F_R\) does not
depend on \(b\).

D4 does not change the worst-case complexity: for generic \(d\) the pointwise
inequality on \(d(b,\cdot)\) vs \(d(b',\cdot)\) fails. It is a correct optional
filter.

\textbf{D5. Bound-based elimination inside the inner min (Bar-Joseph pruning,
dual).} For a \emph{maximization} of similarities, Bar-Joseph et al.\ sort
candidate seams and stop when a remaining candidate cannot beat the current
best given an upper bound on the best unused similarity. The min-cost dual is:
given a current best \(B\) for \(M_v(u,w)\) and a lower bound \(\ell(b)\) on
\(\min_c\bigl(F_R(c,w)+d(b,c)\bigr)\), skip \(b\) when
\(F_L(u,b)+\ell(b) \ge B\). Safe. The supplied code does \textbf{not} implement
it. Worst-case complexity is unchanged (adversarial \(d\) never prunes);
typical complexity can drop, as in the original Bioinformatics experiments.

\subsection{Rules that \emph{look} safe and are not}

\textbf{U1. Scalar dominance across different interfaces.} ``If
\(F_v(u,w) \le F_v(u',w')\) then drop \((u',w')\).'' Unsafe. Different endpoints
have different external options. Counterexample~\ref{ce:5.3}:
\(F_{T_0}(3,4) = 1.5 < 2 = F_{T_0}(2,4)\), but dropping \((2,4)\) loses the
unique global optimum.

\textbf{U2. Same-end, cheaper-start (or same-start, cheaper-end) without a
pointwise condition on \(d\).} ``If \(F_v(u,w) \le F_v(u',w)\) then drop
\((u',w)\).'' Unsafe. The leftward continuation cares about \(u\) vs \(u'\).
In Counterexample~\ref{ce:5.3}, \(F_{T_0}(3,4) < F_{T_0}(2,4)\) with a common
right endpoint \(4\); dropping \((2,4)\) again loses \(1,2,3,4\).

\textbf{U3. Pareto dominance on a scalarization that mixes cost and a single
``centrality'' of the endpoint.} Any scalar that collapses the vector
\(\bigl(d(w,c)\bigr)_{c \notin L(v)}\) can be fooled by putting the only cheap
external neighbour on a coordinate that the scalar underweighted. The safe
version is D4's \emph{pointwise} (Pareto) comparison, which is
high-dimensional and rarely strict.

\textbf{U4. Greedy seam: at each node keep only the orientation whose
\emph{current} child-orders give the cheaper seam.} This is the endpoint-free
algorithm of Section~\ref{sec:fails}. Unsafe by
Counterexample~\ref{ce:5.3}.

\textbf{U5. Reversal identification when \(d \ne d^\top\).} Unsafe by
Theorem~\ref{thm:asym}. Dominating \((w,u)\) by \((u,w)\) discards a genuinely
different directed layout.

\subsection{Resource / bound dominance in a branch-and-bound on flips}

Branch-and-bound is exact whenever its lower bound is \textbf{admissible}: it
must never exceed the minimum possible completion cost of the current partial
choice.

For nonnegative \(d\), zero is an admissible but weak remainder bound. An MST
relaxation is also admissible when constructed on the still-required vertices
(and any required interface terminals), because every feasible Hamiltonian path
completion is a spanning tree after forgetting order and laminar constraints;
the minimum spanning tree can only be cheaper. The OLO constraints make the
feasible completion set smaller, which strengthens the validity of the
relaxation rather than invalidating it.

What an MST bound usually fails to capture is precisely the expensive part of
this problem: required endpoint identities and laminar block compatibility. It
can therefore be weak. Stronger admissible bounds can incorporate interface
constraints, degree relaxations, assignment relaxations, or exact DP values for
unresolved subtrees. None changes the worst-case fact that branch-and-bound may
explore exponentially many flip assignments.

Thus branch-and-bound can be useful on favorable instances, but it is not the
best worst-case formulation: bounded-degree OLO admits a quadratic-time exact
algorithm (Section~\ref{sec:algo}).


\section{Stronger recurrences and alternative formulations}

The existing recurrence is correct (under \(d=d^\top\)) but not computationally
tight. The improvements below are mathematically justified; novelty for its own
sake is avoided.

\subsection{Factored min-plus form --- same exactness, \(O(n^3)\)}

Distribute the min in (R):
\begin{equation}
\label{eq:R3}
M_v(u,w)
\;=\;
\min_{b \in L(L)}
\Bigl(
F_L(u,b)
\;+\;
\min_{c \in L(R)}\bigl(F_R(c,w) + d(b,c)\bigr)
\Bigr).
\end{equation}

\begin{lemma}\label{lem:R3eq}
Over \(\R\cup\{+\infty\}\), (R) and \eqref{eq:R3} are identical.
\end{lemma}

\begin{proof}
\(\min_{b,c} \bigl(A(b) + B(c) + D(b,c)\bigr)
= \min_b \bigl(A(b) + \min_c(B(c)+D(b,c))\bigr)\)
by associativity of \(\min\) and independence of \(A(b)\) from \(c\).
\end{proof}

The inner map \((w,b) \mapsto \min_c\bigl(F_R(c,w)+d(b,c)\bigr)\) does
\textbf{not} depend on \(u\) and is computed once, in time
\(\Theta\bigl(\lambda(L)\,\lambda(R)^2\bigr)\). The outer min is then
\(\Theta\bigl(\lambda(L)^2\,\lambda(R)\bigr)\). Per node this is
\(\Theta\bigl(\lambda(L)\,\lambda(R)\,\lambda(v)\bigr)\) rather than
\(\Theta\bigl(\lambda(L)^2\,\lambda(R)^2\bigr)\).

This is the standard improvement that makes the cba package's
\texttt{order.optimal} an \(O(n^3)\) algorithm; the supplied code implements
the unfactored \(O(n^4)\) quadruple loop.

\eqref{eq:R3} is \textbf{not} a different state. It is the same state with an
equivalent transition, and Theorem~\ref{thm:exactness} transfers verbatim.

In matrix language: if \(A = F_L\), \(B = F_R\), \(D = d|_{L(L)\times L(R)}\),
then \(M_v = A \otimes D \otimes B\) in the min-plus semiring. \eqref{eq:R3} is
parenthesization \((A \otimes D) \otimes B\) vs \(A \otimes (D \otimes B)\),
both \(O(n^3)\) at a balanced root.

\subsection{Directed (asymmetric) exact DP}

Drop the transpose fill. Store both orientations as independent cells:
\[
M_v^{\rightarrow}(u,w),\qquad
M_v^{\leftarrow}(w,u),
\]
and use the true \(F_L(u,b) = M_L^{\text{oriented}}(u,b)\) without identifying
reversals. Transitions stay (R) / \eqref{eq:R3}, with \(d(b,c)\) \emph{not}
replaced by \(d(c,b)\). Exact for arbitrary real \(d\), at a constant-factor
increase in stored orientation states and the same asymptotic time. This claim
follows from the same endpoint-conditioned induction as
Theorem~\ref{thm:exactness}, now without identifying reverse states.

This is the minimal change that restores Theorem~\ref{thm:exactness} without
the hypothesis \(d = d^\top\).

\subsection{Stronger general result: an asymptotically optimal \(O(n^2)\)
exact algorithm}
\label{sec:algo}

The cubic endpoint-table DP is \textbf{not} asymptotically optimal for
arbitrary binary trees.

Brandes (2007), while reviewing prior work, states that for a tree of maximum
node degree \(d\), Deineko and Tiskin gave an exact algorithm with
\[
O(4^d n^2) \text{ time}, \qquad O(4^d n) \text{ additional space}.
\]
For binary trees, \(d\) is a fixed constant, hence this becomes
\[
O(n^2) \text{ time}, \qquad O(n) \text{ additional space}.
\]
This matches the natural \(\Omega(n^2)\) lower bound in the
explicit-distance-matrix model.

\begin{theorem}[quadratic input-query lower bound]\label{thm:quadlb}
In a model where the entries of an arbitrary symmetric dissimilarity matrix are
supplied explicitly or by random-access queries, every exact algorithm has
worst-case \(\Omega(n^2)\) input/query complexity.
\end{theorem}

\begin{proof}
Run the algorithm on the all-zero symmetric matrix. Suppose some off-diagonal
pair \(\{i,j\}\) is never inspected. Let \(v=\LCA(i,j)\). By orienting the child
of \(v\) containing \(i\) so that \(i\) is its rightmost leaf, orienting the
other child so that \(j\) is its leftmost leaf, and choosing arbitrary feasible
orientations elsewhere, there exists a tree-consistent order in which \(i,j\)
are adjacent. Now form a second input identical on all queried entries but set
\(d(i,j)=d(j,i)=-1\). The optimum of the all-zero input is \(0\), while the
modified input has a feasible order of cost at most \(-1\). The algorithm sees
the same answers to every query on both inputs and therefore cannot be correct
on both. Hence in the worst case every unordered pair must be inspected, giving
\(\Omega(\binom{n}{2})=\Omega(n^2)\).
\end{proof}

The quadratic algorithm comes from the PQ-tree / minimum-weight tree-
shortcutting viewpoint rather than from storing every endpoint table and
performing a min-plus product at each internal node. Its core uses
\textbf{sweeping paths} and specialized layered acyclic graphs. For bounded
degree, the number of child-subset configurations is constant; a global counting
argument bounds the relevant triples/quadruples by \(O(n^2)\), and shortest
paths in the layered DAGs are computed in time proportional to their arcs. This
is a genuine change of decomposition, not merely a faster implementation of
\eqref{eq:R3}.

The quadratic result is cited here through Brandes specifically as an exact
algorithmic result for tree-compatible \textbf{leaf ordering}. Therefore no
auxiliary reduction from this linear-path OLO objective to a cyclic TSP
objective is needed for the complexity conclusion. This avoids importing extra
representability assumptions about a dummy terminal into the fixed-tree
feasible family.

\textbf{Consequence.} The theoretically best asymptotic target on binary trees
is quadratic, not cubic. The endpoint-table recurrence remains attractive
because it is considerably simpler to prove and implement, but it is not state
of the art in worst-case complexity.

\subsection{Brandes (2007): practical \(O(n^2\log n)\) for complete /
near-complete trees}

For complete binary trees, and more generally trees of height
\(\log_2 n+O(1)\) after controlled completion, Brandes gives a particularly
simple algorithm using canonical bit labels and position-based predecessor
sets. It runs in \(O(n^2\log n)\) time and \(O(n)\) additional space.

Its key positional lemma fails for incomplete trees: fixing a leaf at a position
no longer uniquely determines the partition of the remaining leaves. Padding a
height-\(h\) tree to completeness may create \(2^h\) leaves; for a comb this is
exponential. Thus Brandes is a practical specialization, not the general
quadratic method of Section~\ref{sec:algo}.

\subsection{Comb (caterpillar) specialization: \(O(n^2)\)}
\label{sec:comb}

If \(T_k=(T_{k-1},k)\), leaf \(k\) is an endpoint of every arrangement of
\(T_k\). The only nontrivial state is which older leaf forms the opposite
endpoint. There are \(O(k)\) such states at layer \(k\), and the transition can
be restricted to the finite support of the preceding comb table, yielding
\(O(n^2)\) total time.

The supplied code still spends \(\Theta(n^3)\) on a comb because it scans many
infinite entries in dense full matrices. Iterating only over finite support
recovers the quadratic behavior on this shape.

\subsection{Alternative formulations: what applies and what does not}

\textbf{Shortest path / DAG.} A shortest-path interpretation \emph{does} help,
but not as ``run Dijkstra on the leaves.'' The useful graph is a layered
\textbf{state graph} of sweeping configurations, as in the quadratic
bounded-degree algorithm. Because it is acyclic, topological shortest-path DP is
preferable to Dijkstra and also works with negative edge weights.

The endpoint recurrence (R) can likewise be viewed as a shortest hyperpath in
an acyclic state hypergraph. That viewpoint alone does not reduce complexity;
the improvement comes from choosing a state graph with only \(O(n^2)\) relevant
transitions for bounded degree.

\textbf{PQ-tree formulation.} This is directly relevant. A rooted binary tree
that allows only child swaps is a restricted PQ-tree. Classical PQ-tree TSP
machinery gives a polynomial exact formulation, and the later
tree-shortcutting specialization is what reaches quadratic time for bounded
degree.

\textbf{Min-cost flow.} No comparably direct compact network-flow formulation
is needed to obtain the optimal asymptotic bound. Encoding ``visit each leaf
once'' and laminar contiguity through a generic flow model tends to reintroduce
TSP-like degree/order structure. Since \(\Omega(n^2)\) time is already forced
by an explicit arbitrary cost matrix, a flow formulation could at best match,
not beat, the input-reading bound.

\textbf{Held--Karp / meet-in-the-middle.} These are exact but waste the laminar
tree structure. Held--Karp ranges over arbitrary subsets and is exponential;
meet-in-the-middle over the \(n-1\) flip bits is also exponential. Both are
inferior to the tree-specific polynomial algorithms.

\textbf{Knuth/Yao, divide-and-conquer optimization, SMAWK/Monge.} These require
additional monotonicity or quadrangle/Monge structure on a \emph{known}
indexing. The leaf order here is itself the decision variable, and a generic
Euclidean trait-distance matrix is not Monge in the arbitrary tip labelling.
They may apply to special structured distance families, but not to the stated
general problem without extra assumptions.

\textbf{Circular OLO.} Adding a wrap-around term gives a different objective,
but can be handled by the endpoint machinery or the tour/PQ-tree formulation.
It is not an acceleration of the linear objective.

\subsection{How close is each formulation to optimal?}

\begin{itemize}
\item supplied literal R recurrence: worst-case \(\Theta(n^4)\);
\item factored endpoint DP \eqref{eq:R3}: \(\Theta(n^3)\) worst-case;
\item Brandes on complete/near-complete trees: \(O(n^2\log n)\), linear
  additional space;
\item bounded-degree tree-shortcutting/PQ-tree algorithm: \(O(n^2)\) time and
  \(O(n)\) additional working space for binary trees;
\item explicit arbitrary \(n\times n\) dissimilarity input: \(\Omega(n^2)\) time
  in the standard RAM/query model.
\end{itemize}

Hence the general bounded-degree quadratic algorithm is asymptotically optimal
with respect to explicit-input size. There is no cubic barrier.

\section{Complexity}
\label{sec:complexity}

\subsection{The implementation as written (R)}

Let \(a_v=\lambda(L)\) and \(b_v=\lambda(R)\) for the two children of an
internal node \(v\). Let \(q_L\) denote the number of \textbf{finite ordered
endpoint pairs} in the full child table \(F_L\). For a leaf child, \(q_L=1\).
If \(L\) is internal with grandchildren of sizes \(p\) and \(q\), then
\(q_L=2pq\), because exactly the cross-grandchild ordered pairs are feasible.

The literal R implementation is not a completely dense four-loop evaluation,
because it tests \texttt{is.finite(Lcost)} before constructing the
length-\(b_v\) vector over \(c\). Its work at node \(v\) separates into:

\begin{enumerate}[label=\arabic*.]
\item \(\Theta(a_v^2 b_v)\) scalar \texttt{(u,w,b)} checks;
\item \(\Theta(a_v b_v^2)\) extraction/scanning of the length-\(b_v\)
  \texttt{Rcol} vectors across all \((u,w)\) pairs (up to constant and R-copy
  semantics);
\item \(\Theta(q_L b_v^2)\) arithmetic and \texttt{which.min} work for the
  finite \((u,b)\) pairs, because each such pair is revisited for every \(w\)
  and scans all \(c\in L(R)\).
\end{enumerate}

Thus a support-aware arithmetic model is
\[
T_v^{\text{literal}}
=
\Theta\bigl(a_v^2 b_v + a_v b_v^2 + q_L b_v^2\bigr),
\]
up to implementation-level copying constants. A completely dense evaluation of
recurrence (R), ignoring the \texttt{is.finite} shortcut, would instead cost
\(\Theta(a_v^2b_v^2)\).

The distinction does \textbf{not} change the worst case. At a balanced node
whose left child is itself balanced, \(q_L=\Theta(a_v^2)\) and
\((a_v,b_v)=\Theta(s_v)\), so that node costs \(\Theta(s_v^4)\). Hence a
perfectly balanced tree satisfies
\[
T(n)=2T(n/2)+\Theta(n^4)=\Theta(n^4).
\]
Therefore the supplied implementation has \textbf{worst-case \(\Theta(n^4)\)
time}.

For a comb, the finite support of the large child is only linear but the
\texttt{(u,w,b)} checks are still quadratic at layer \(s\), giving
\(\Theta(s^2)\) per layer and \(\Theta(n^3)\) total. Mixed tree shapes can be
strictly cheaper than the dense formula \(\sum_v a_v^2b_v^2\), so that formula
should be read as a dense-recurrence bound, \textbf{not} as an exact count for
every execution of this R code.

\textbf{Space.} The number of stored rectangular DP cells is
\[
\sum_v a_vb_v=\binom{n}{2},
\]
because each unordered pair of leaves is separated for the first time at its
unique LCA. Hence all DP tables and traceback pointers use \(\Theta(n^2)\)
words. The temporary full child matrices can also reach \(\Theta(n^2)\), so
total working memory is \(\Theta(n^2)\).

\subsection{Factored transition \eqref{eq:R3}}

Per node \(\Theta(a_v b_v s_v)\). Using \(s_v \le n\) and the identity above,
\[
T_{\mathrm{(R3)}}(T)
\;=\;
\Theta\Bigl(\sum_v a_v b_v s_v\Bigr)
\;\le\;
n \sum_v a_v b_v
\;=\;
O(n^3).
\]
The bound is tight: the root of a balanced tree already contributes
\(\Theta(n^3)\). The same \(\Theta(n^3)\) holds for a comb
(\(\sum_s s\cdot 1\cdot s = \Theta(n^3)\)) unless finite-entry iteration or
the comb specialization is used.

Space is unchanged: \(\Theta(n^2)\).

\subsection{Special cases, pruning, typical}

\begin{center}
\renewcommand{\arraystretch}{1.4}
\begin{tabular}{@{}l l l@{}}
\toprule
Setting & Time & Notes \\
\midrule
(R), worst case & \(\Theta(n^4)\) & balanced tree, generic \(d\) \\
(R), comb & \(\Theta(n^3)\) & R loops; \(\Theta(n^2)\) if only finite
\(F\)-entries drive the \(b\)-loop \\
(R3), any binary tree & \(\Theta(n^3)\) & tight for this recurrence at a
balanced root \\
(R3), comb, finite-entry loops & \(\Theta(n^2)\) & equivalent to
Section~\ref{sec:comb} \\
Brandes, complete / height \(\log n+O(1)\) & \(O(n^2 \log n)\) & \(O(n)\)
additional space \\
Deineko--Tiskin bounded-degree tree shortcutting & \(O(n^2)\) for binary trees
& asymptotically optimal; \(O(n)\) additional space \\
D4 / D5 pruning & still \(\Theta(n^3)\) worst case & typical-case only \\
Endpoint-free greedy & \(\Theta(n)\) after child orders, but \textbf{inexact}
& not a competitor \\
Brute force on \(\phi\) & \(\Theta(n\,2^{n-1})\) & used as an oracle, not an
algorithm \\
\bottomrule
\end{tabular}
\end{center}

\textbf{Where the objective adds a factor.} The
\(\lambda(L)\lambda(R)\) interfaces are forced by Theorem~\ref{thm:nodrop}. The
remaining \(\times \lambda(v)\) in \eqref{eq:R3} is the cost of minimizing over
the seam against an \emph{arbitrary} \(d\). If \(d\) were e.g.\ rank-\(1\) or a
1-D Euclidean matrix in a \emph{known} order, that factor could drop; under the
problem as stated it is inherent to the transition.

\subsection{The \(n = 500\) driver}

The comments advertise a 500-tip example. Under (R) a balanced 500-leaf tree is
\(\sim 500^4 / 14 \approx 4 \times 10^9\) inner operations, nested in R
interpreter loops --- not a realistic run. Under \eqref{eq:R3} it is
\(\sim 500^3 / 4 \approx 3 \times 10^7\) at the root plus a geometric series,
which is routine in C and uncomfortable in pure R. A quadratic bounded-degree
implementation is theoretically better still, but much more complex to
engineer. The exactness theorem does not depend on \(n\); tractability does.

\section{Exhaustive search as a validation tool}

\subsection{What brute force can establish}

Enumeration of all \(2^{n-1}\) flip vectors, for a \emph{fixed} \((T,d)\),
computes the true \(\min J\) on that instance. Agreement of a DP with that value
on a collection of instances is \textbf{evidence}, not a proof of exactness for
all \((T,d)\).

It \emph{can} establish:

\begin{itemize}
\item existence of counterexamples (one failing instance is a proof of
  inexactness --- Theorems~\ref{thm:asym} and~\ref{thm:noendpoint} were obtained
  this way);
\item that two recurrences (R) and \eqref{eq:R3} agree on the tested set;
\item that a proposed dominance rule is unsafe (a dominated state used by the
  unique optimum);
\item that an implementation bug exists on small \(n\).
\end{itemize}

\subsection{What brute force cannot establish}

\begin{itemize}
\item Exactness for all \(n\). A bug that first appears at \(n = 11\) (or on a
  tree shape not sampled, or on a measure-zero algebraic variety of \(d\))
  will be missed. The authors' report of agreement for \(n = 4..10\) is
  consistent with Theorem~\ref{thm:exactness} but is not that theorem.
\item Exactness for a reduced state. Agreement after dropping a cell on random
  \(d\) may mean only that the dropped cell was not needed \emph{on those
  \(d\)}.
\item Asymptotic complexity.
\end{itemize}

\subsection{What \emph{would} constitute a proof of equivalence between (R) and
brute force}

A proof that, for every \((T,d)\) in the stated class,
\[
\min M_\rho(T,d)
\;=\;
\min_{\phi \in \{0,1\}^{n-1}} J(\pi_\phi; d),
\]
together with a proof that traceback emits a \(\pi_\phi\) realizing the left
side. That is Theorem~\ref{thm:exactness}. It is an induction on \(T\), not a
finite check.

A finite check \emph{plus} a proof that ``if it works for all trees with
\(n \le n_0\) then it works for all \(n\)'' would require a reduction of large
instances to small ones. No such reduction exists for this DP: the inductive
step \emph{is} the proof, and it does not factor through a bounded \(n\).

\subsection{A systematic counterexample methodology}

The search that produced Sections~\ref{sec:exact} and~\ref{sec:minimalce}, and
that failed to break Theorem~\ref{thm:exactness}, is:

\begin{enumerate}[label=\arabic*.]
\item \textbf{Independent oracle:} enumerate \(\pi_\phi\), compute \(J\) in
  floating point, take the min. Do not reuse DP code in the oracle.
\item \textbf{Instance families:} (a) all binary tree shapes for \(n \le 6\);
  (b) random recursive binary trees; (c) combs; (d) perfectly balanced; (e) the
  structured four-leaf tree of Counterexample~\ref{ce:5.3}.
\item \textbf{Matrices:} Euclidean traits of dimension \(1,2,5,10\); generic
  symmetric; generic asymmetric; all-zero; all-off-diagonal-equal; symmetric
  with negative entries; quantized (ties).
\item \textbf{Predators, not just agreement tests:}
  \begin{itemize}
  \item endpoint-free greedy (must fail if endpoints are necessary);
  \item (R) with transpose fill on asymmetric \(d\) (must fail if symmetry is
    necessary);
  \item \eqref{eq:R3} vs (R) (must agree if the identity is implemented);
  \item directed two-orientation DP vs brute (must agree if the directed DP is
    right).
  \end{itemize}
\item \textbf{Tie and \(\mathrm{NaN}\) probes.} Quantized \(d\); constant trait
  columns (the \texttt{scale} hazard).
\item \textbf{Cost-vs-order consistency check:} the reported scalar must equal
  \(J\) of the reconstructed permutation, independently computed. This is what
  caught the ``ghost cost'' in Counterexample~\ref{ce:3.4}'s \(n=4\) case.
\end{enumerate}

A falsification of Theorem~\ref{thm:exactness} would be any instance in class
(rooted binary \(T\), symmetric finite \(d\)) on which step 6 fails or the DP
min differs from the oracle. None was found; that is evidence, and the proof is
Section~\ref{sec:exact}.


\section{Counterexample search}

Adversarial targets and outcomes.

\subsection{Incorrect state compression}

Endpoint-free compression is inexact: Counterexample~\ref{ce:5.3}, \(n=4\), gap
\(101.5\) vs \(2\). The counterexample is already a theorem-level falsification;
no failure frequency is needed for the conclusion.

Dropping one endpoint while keeping the other is inexact by
Theorem~\ref{thm:twoend}; the same structured instance is a witness.

\subsection{Missing history}

Any history beyond \((v,u,w)\) is redundant
(Theorem~\ref{thm:suff}). Search for a counterexample to \emph{sufficiency} of
that state: none exists under the exchange argument, which does not depend on
\(n\).

\subsection{Invalid optimal-substructure assumption}

The unconstrained-subtree version fails
(Counterexample~\ref{ce:5.3}). The endpoint-sliced version holds
(Theorem~\ref{thm:optend}). The implementation uses the latter.

\subsection{Invalid dominance}

U1--U5 of Section~\ref{sec:dom} are each refuted by
Counterexample~\ref{ce:5.3} or~\ref{ce:3.4}.

\subsection{Incorrect transition ordering}

At the node level, any post-order is a valid evaluation order because every
parent depends only on its descendants. At the finer endpoint-state level the
dependency structure is an acyclic hypergraph. The recursive implementation
respects this partial order automatically, so there is no transition that uses
a child table before that child has been solved.

\subsection{Boundary cases, ties, degeneracy}

\begin{center}
\renewcommand{\arraystretch}{1.4}
\begin{tabular}{@{}l p{10.5cm}@{}}
\toprule
Probe & Result under (R) with symmetric \(d\) \\
\midrule
\(n=2\) (single cherry) & exact: the only cost is \(d(1,2)\) \\
\(n=3\) & exact under symmetry; \textbf{inexact} without
(Counterexample~\ref{ce:3.4}) \\
all-zero \(d\) & exact, cost \(0\), every order optimal \\
equal off-diagonal \(d \equiv 1\) & exact, cost \(n-1\) \\
negative symmetric \(d\) & exact (non-negativity unused) \\
quantized \(d\) (many ties) & exact cost; arbitrary choice among optima \\
comb and balanced, \(n \le 8\) & exact \\
random Euclidean/symmetric small instances & empirical agreement with brute
force; evidence only \\
constant trait \(\Rightarrow\) \texttt{NaN} after \texttt{scale} & \textbf{not
exact: IEEE \texttt{NaN} poisons the min} (implementation hazard, outside
Theorem~\ref{thm:exactness}'s finite-\(d\) hypothesis) \\
\bottomrule
\end{tabular}
\end{center}

\subsection{Asymmetric / directed inputs}

Random directed tests are useful only as falsification aids. The theorem-level
smallest witness is Counterexample~\ref{ce:3.4} with \(n=2\); a four-leaf
balanced tree can additionally expose cost/order inconsistency through
transpose-filled child states.

\subsection{Locally inferior partial solutions that are globally superior}

This is the \emph{content} of Counterexample~\ref{ce:5.3}, and is exactly why
the state must retain currently-suboptimal interfaces. The implementation does
retain them (all cross-child pairs). The endpoint-free algorithm does not.

\subsection{Multiple histories, same nominal state, different futures}

Under the state \((v,u,w)\), Theorem~\ref{thm:suff} says this cannot happen.
Under a compressed state it does happen, and~\ref{ce:5.3} is the picture.

\subsection{What was \emph{not} found}

No counterexample to Theorem~\ref{thm:exactness} on the class it claims. That
is \textbf{not} a proof; the proof is Section~\ref{sec:exact}. The search is
reported here because the mandate was to try to break the algorithm
\emph{before} proving it, and the breakage that exists (asymmetric \(d\),
endpoint-free compression, \texttt{NaN}) was found by that search, not by
accepting the recurrence.

\section{Empirical evidence versus proof}

\begin{center}
\renewcommand{\arraystretch}{1.35}
\begin{tabular}{@{}p{11.2cm} p{4.3cm}@{}}
\toprule
Claim & Status \\
\midrule
(R) is exact for rooted binary \(T\) and symmetric finite \(d\) &
\textbf{Proven} (Theorem~\ref{thm:exactness}) \\
(R3) is equivalent to (R) & \textbf{Proven} (Lemma~\ref{lem:R3eq}) \\
\((v,u,w)\) is sufficient & \textbf{Proven} (Theorem~\ref{thm:suff}) \\
neither endpoint is droppable & \textbf{Proven}
(Theorem~\ref{thm:twoend}, Counterexample~\ref{ce:5.3}) \\
unconstrained-subtree optimal substructure fails & \textbf{Proven}
(Theorem~\ref{thm:noendpoint}, Counterexample~\ref{ce:5.3}) \\
(R) is inexact for asymmetric \(d\) & \textbf{Proven}
(Theorem~\ref{thm:asym}, Counterexample~\ref{ce:3.4}) \\
D1--D5 are safe dominance rules & \textbf{Proven}
(Section~\ref{sec:dom}) \\
U1--U5 are unsafe & \textbf{Proven} by counterexample
(Section~\ref{sec:dom}) \\
\(\sum_v a_v b_v = \binom{n}{2}\) & \textbf{Proven} (LCA counting) \\
\(T_{\mathrm{(R)}} = \Theta(n^4)\) worst case; \(T_{\mathrm{(R3)}} =
\Theta(n^3)\) & \textbf{Proven} (Section~\ref{sec:complexity}) \\
Brandes' position-based algorithm does not extend directly to arbitrary-height
trees & \textbf{Proven in Brandes' analysis} \\
a quadratic exact algorithm exists for bounded-degree trees &
\textbf{Established in the literature} (Deineko--Tiskin, cited by Brandes) \\
independent small-instance brute force agrees with (R) on tested symmetric
cases & \textbf{Empirical only} --- supports, does not replace,
Theorem~\ref{thm:exactness} \\
the code comment's claimed \(n=4..10\) brute-force agreement &
\textbf{Empirical claim in the supplied material only} \\
D4/D5 help on typical trait data & \textbf{Conjectured / historically
reported} (Bar-Joseph 2001 tables); not re-measured here \\
\texttt{rotateConstr} realizes the DP order on every ape tree &
\textbf{Not proven here}; outside the DP. A mismatch would be an ape bug or a
contract bug, not a recurrence bug \\
Euclidean trait \(d\) has extra structure that a specialized OLO could exploit
& \textbf{Conjectured}; no algorithm given \\
\bottomrule
\end{tabular}
\end{center}

No empirical row is promoted to a theorem in Section~\ref{sec:verdict}.

\section{Final theoretical verdict}
\label{sec:verdict}

\subsection{Exactness}

For the intended domain --- a rooted strictly binary tree and a finite symmetric
matrix \(d\) --- the algorithm is \textbf{globally exact}
(Theorem~\ref{thm:exactness}). Symmetry is a clean sufficient condition because
it makes every subtree optimum table reversal-invariant.

This is \textbf{not an instance-wise ``if and only if'' characterization}. Some
asymmetric matrices can accidentally satisfy the needed reversal invariance.
What is true is that arbitrary asymmetric input is outside the guaranteed
correct domain: the smallest counterexample has \(n=2\)
(Theorem~\ref{thm:asym}), and larger instances can additionally create reported
costs inconsistent with the reconstructed directed order.

\textbf{Not defined} for polytomies or unrooted trees; the implementation
correctly rejects the former and silently assumes the latter are rooted.

Non-negativity, the triangle inequality, and the trait-Euclidean origin of \(d\)
are \textbf{not} required for exactness.

\subsection{Required state}

For the standard \textbf{subtree-composition} DP, the minimal generic interface
is the two endpoint identities. Under symmetric \(d\), the state index can be
the unordered pair \(\{u,w\}\); under asymmetric \(d\), it must be the ordered
pair \((u,w)\). The subtree node is derivable as \(v=\LCA(u,w)\), and
\(F_v(u,w)\) is the value stored at the state rather than an additional state
coordinate.

Both endpoint identities are necessary for arbitrary continuation contexts
(Theorem~\ref{thm:twoend}); the full internal order and flip history are not.
Theorem~\ref{thm:nodrop} shows that no generic endpoint cell can be pruned
uniformly from this compositional DP, although a different exact algorithm need
not materialize all such cells simultaneously.

\subsection{Optimal substructure}

Holds for the endpoint-conditioned state (Theorem~\ref{thm:optend}) and fails
for the endpoint-free state (Theorem~\ref{thm:noendpoint}). There is no
greedy-choice property. Child endpoint values are reused across many parent
transitions; the state-dependency structure is an acyclic hypergraph induced by
the rooted tree, not merely the node tree itself.

\subsection{Dominance}

Safe to eliminate: duplicate reverse cells when \(d = d^\top\) (D2); infeasible
same-child pairs (D3); same-interface worse internals (D1);
pointwise-seam-dominated candidates (D4); bound-pruned candidates in the inner
min (D5).

Unsafe to eliminate: any cell solely because it has worse internal cost than a
cell with a \emph{different} interface (U1, U2, Counterexample~\ref{ce:5.3}).

\subsection{Complexity}

\begin{center}
\renewcommand{\arraystretch}{1.35}
\begin{tabular}{@{}l l l l@{}}
\toprule
Algorithm & Time & Additional/working space & Domain \\
\midrule
supplied literal R code & \(\Theta(n^4)\) worst case & \(\Theta(n^2)\) &
general binary, symmetric finite \(d\) \\
factored endpoint DP (R3) & \(\Theta(n^3)\) worst case & \(\Theta(n^2)\) & same
\\
support-aware comb DP & \(\Theta(n^2)\) & shape-dependent & combs \\
Brandes 2007 & \(O(n^2\log n)\) & \(O(n)\) & complete / height
\(\log n+O(1)\) \\
Deineko--Tiskin bounded-degree algorithm & \(O(n^2)\) for binary trees &
\(O(n)\) & general bounded-degree trees \\
brute force & \(\Theta(n2^{n-1})\) & \(O(n)\) besides input & validation oracle
\\
\bottomrule
\end{tabular}
\end{center}

The tightest justified asymptotic bound for the \textbf{general binary-tree
problem} with explicit pairwise costs is \(\Theta(n^2)\) time: the literature
gives an \(O(n^2)\) exact bounded-degree algorithm, while reading/querying an
arbitrary dense dissimilarity matrix requires \(\Omega(n^2)\) in the standard
model. The algorithm's extra working memory can be linear, although the
supplied dense matrix \(d\) itself already occupies \(\Theta(n^2)\) storage.

\subsection{Potential improvements (theoretically justified,
exactness-preserving)}

\begin{enumerate}[label=\arabic*.]
\item \textbf{Immediate implementation change: replace (R) by (R3).} Same proof
  and state semantics, worst case \(O(n^4)\to O(n^3)\); this is the simplest
  high-impact improvement to the supplied R code.
\item \textbf{Exploit finite support.} Do not materialize or scan \texttt{Inf}
  regions of child endpoint matrices. This gives strong shape-dependent savings
  and quadratic behavior on comb-like trees.
\item \textbf{For asymptotic optimality, change decomposition.} Implement the
  bounded-degree PQ-tree / tree-shortcutting algorithm: \(O(n^2)\) time and
  \(O(n)\) additional space for binary trees. It is mathematically stronger but
  much more complex than (R3).
\item \textbf{Use Brandes for complete or near-complete trees.} It is simpler
  than the general quadratic method and uses linear additional memory.
\item \textbf{Use safe dominance/bounds only as secondary accelerators.}
  D4/D5 and admissible branch-and-bound bounds may help typical instances but
  do not improve the endpoint DP's worst-case asymptotic guarantee.
\item \textbf{If directed costs are needed, store both endpoint orientations}
  rather than using the transpose shortcut.
\item \textbf{Guard \texttt{scale()} against zero-variance traits} so the
  distance matrix stays finite.
\end{enumerate}

Do \textbf{not} drop an endpoint, keep only one order per subtree, or identify
reversals for directed costs. Generic Dijkstra-on-leaves, Held--Karp,
meet-in-the-middle, and standard Monge/Knuth optimizations do not exploit the
right structure.

\subsection{Remaining gaps}

\begin{itemize}
\item No cubic lower bound applies to the problem: Section~\ref{sec:algo}
  already records an exact \(O(n^2)\) bounded-degree algorithm. What remains
  open \emph{within the simple endpoint-table/min-plus decomposition used by
  (R3)} is how far one can reduce its work without changing decomposition or
  exploiting additional structure in \(d\). The document therefore does
  \textbf{not} claim that \(\Theta(n^3)\) is an intrinsic barrier for OLO.
\item \texttt{ape::rotateConstr} is not part of the proof. Exactness is claimed
  for \texttt{olo\_phylo}'s \texttt{(order, cost)}, not for the plotted
  embedding.
\item Numerical stability of the \texttt{which(table == min\_cost)} tie-break
  under floating-point noise is not analysed; it does not affect the
  \emph{value} of the min, only which optimum is selected.
\item Euclidean-trait structure (low-dimensional embeddings,
  identity-minus-correlation matrices, etc.) is unused. Whether it yields a
  faster \emph{exact} algorithm is open.
\end{itemize}

\subsection{Confidence}

\textbf{Theorem-level:} exactness of (R) and (R3) on rooted binary trees with
symmetric finite \(d\); lack of a correctness guarantee for arbitrary
asymmetric \(d\); sufficiency and necessity of the two-endpoint state; failure
of endpoint-free optimal substructure; safe/unsafe dominance classification
D1--D5 / U1--U5; \(\Theta(n^4)\) vs \(\Theta(n^3)\) arithmetic complexity.

\textbf{Empirical:} supplementary small-instance brute-force checks are
consistent with Theorem~\ref{thm:exactness} on symmetric inputs and readily
produce failures on directed inputs. These checks are validation aids only;
none of the theorem-level claims relies on their success.

\textbf{Literature-level conclusion:} the general binary-tree problem admits an
asymptotically optimal quadratic-time exact algorithm for explicit pairwise
cost input. The cubic endpoint DP is simpler, but not asymptotically optimal.

\section*{References and external cross-checks}

\begin{enumerate}[label=\arabic*.]
\item Z.\ Bar-Joseph, D.~K. Gifford, and T.~S. Jaakkola, ``Fast optimal leaf
  ordering for hierarchical clustering,'' \emph{Bioinformatics} 17(Suppl.\ 1),
  S22--S29, 2001. DOI: 10.1093/bioinformatics/17.suppl\_1.S22. The paper
  presents the original practical OLO framework and an \(O(n^4)\) basic
  algorithm with further improvements.
\item R.~E. Burkard, V.~G. Deineko, and G.~J. Woeginger, ``The Travelling
  Salesman and the PQ-Tree,'' \emph{Mathematics of Operations Research} 23(3),
  613--623, 1998; erratum 24(1), 262--272, 1999. It gives an \(O(2^d n^3)\)
  exact algorithm for bounded-degree PQ-tree-restricted TSP.
\item V.~G. Deineko and A.~Tiskin, ``Fast minimum-weight double-tree
  shortcutting for Metric TSP,'' WEA 2007 / later JEA version. The
  bounded-degree sweeping-path algorithm runs in \(O(4^d n^2)\) time and
  \(O(4^d n)\) working space; for binary trees this is quadratic time and
  linear additional space.
\item U. Brandes, ``Optimal leaf ordering of complete binary trees,''
  \emph{Journal of Discrete Algorithms} 5(3), 546--552, 2007. DOI:
  10.1016/j.jda.2006.09.003. Brandes recapitulates the binary-tree endpoint DP
  as \(\Theta(n^3)\), cites the quadratic Deineko--Tiskin result, and gives a
  simple \(O(n^2\log n)\)-time, \(O(n)\)-space method for complete /
  near-complete trees.
\item \texttt{cba::order.optimal} documentation states \(O(n^3)\) time for its
  binary-tree OLO implementation.
\end{enumerate}

These references support the historical and asymptotic context. In particular,
Brandes (2007) explicitly states that tree-compatible ordering with bounded
node degree is polynomial and cites the Deineko--Tiskin
\(O(4^d n^2)\)-time, \(O(4^d n)\)-space exact result; for fixed-degree binary
trees this is quadratic time. The proof of exactness of the supplied recurrence
under symmetric finite \(d\), the endpoint-sufficiency exchange argument, and
the counterexamples in this document are self-contained.

\end{document}
